\section{Methodology}

\subsection{Gradient Abnormality Detection Pipeline}

We extend GAIA \cite{chen2023exploring} from out-of-distribution (OOD) detection to minority group detection within in-distribution data. The pipeline operates in three stages:

\textbf{Stage 1: GradCAM Extraction}
For test sample $x_i$ with predicted class $c$, compute gradient of class score $S_c$ w.r.t.~final convolutional layer activations $A^k$ (channel $k$):
\begin{equation}
\alpha_c^k = \frac{1}{Z} \sum_{i,j} \frac{\partial S_c}{\partial A_{i,j}^k}
\end{equation}
where $Z$ is the spatial dimension. Gradient-weighted activation map:
\begin{equation}
L_{\text{GradCAM}} = \text{ReLU}\left(\sum_k \alpha_c^k A^k\right)
\end{equation}

\textbf{Stage 2: GAIA-Z Computation (Zero-Deflation)}
Measure gradient sparsity via near-zero element ratio:
\begin{equation}
\text{GAIA-Z}(x) = \frac{|\{g_{ij} : |g_{ij}| < \epsilon\}|}{|G|}
\end{equation}
where $G = \nabla_A S_c$ is the gradient tensor, $\epsilon = 10^{-5}$ is the near-zero threshold. Higher GAIA-Z indicates denser gradients (more abnormal). Minority samples expected to have lower GAIA-Z (fewer near-zero gradients) due to gradient scattering.

\textbf{Stage 3: Statistical Testing}
Split test samples by group membership: $\mathcal{M}$ (minority), $\mathcal{J}$ (majority). Compute:
\begin{itemize}
\item \textbf{Divergence:} $\Delta = |\text{mean}(\text{GAIA-Z}_\mathcal{M}) - \text{mean}(\text{GAIA-Z}_\mathcal{J})|$
\item \textbf{Significance:} Welch's t-test (unequal variance), $H_0: \mu_\mathcal{M} = \mu_\mathcal{J}$
\item \textbf{Effect size:} Cohen's $d = \frac{\mu_\mathcal{M} - \mu_\mathcal{J}}{\sqrt{\frac{\sigma_\mathcal{M}^2 + \sigma_\mathcal{J}^2}{2}}}$
\end{itemize}

\textbf{Detection criterion:} Divergence $\geq$0.2, $p<0.01$, Cohen's $d\geq$0.8 (large effect).

\subsection{Causal Mechanism: 4-Step Chain}

Our hypothesis proposes gradient abnormality arises causally from spurious conflict:

\textbf{Step 1:} Model learns spurious shortcut (e.g., 90\% background-class correlation in training).
\textbf{Step 2:} Minority samples create conflict --- spurious feature (waterbird on land) contradicts expected shortcut (waterbird on water).
\textbf{Step 3:} Conflict manifests as gradient scattering --- attribution attempts to explain prediction using both spurious region (background) and core region (bird), producing dense, noisy gradients.
\textbf{Step 4:} Higher GAIA-Z score detects scattering, enabling minority group identification.

\textbf{Falsifiable prediction (Step 2 validation):} If spurious mismatch \emph{causes} abnormality, then background swap (minority $\rightarrow$ majority background) should reduce GAIA-Z by $\geq$30\%.

\subsection{Background Augmentation Causality Test}

To validate Step 2 (conflict $\rightarrow$ scattering), we perform an intervention:

\textbf{Procedure:}
\begin{enumerate}
\item Select minority samples $\{x_i\}_{i \in \mathcal{M}}$ with mismatched backgrounds
\item Apply background swap: $\tilde{x}_i = \text{SegFormer}(x_i, \text{bg}_{\text{majority}})$
   (Replace minority background with majority-group background using semantic segmentation)
\item Compute GAIA-Z before ($z_i$) and after ($\tilde{z}_i$) augmentation
\item Paired t-test: $H_0: \text{mean}(z - \tilde{z}) = 0$
\end{enumerate}

\textbf{Success criterion:} Mean reduction $\geq$30\%, $p<0.05$, Cohen's $d\geq$0.5 (medium effect).

\textbf{Rationale:} Causal intervention removes spurious conflict $\rightarrow$ gradient scattering should decrease. Non-causal explanation (e.g., sample complexity) would not change after background swap.

\subsection{Spatial Gradient Regularization}

\subsubsection{Algorithm Overview}

We propose a training-time regularization that penalizes gradients in spurious regions identified via GradCAM difference maps:

\textbf{Input:} Training batch $(x, y)$, model $f_\theta$
\textbf{Output:} Loss $L = L_{\text{CE}} + \lambda \cdot L_{\text{spatial}}$

\textbf{Step 1: GradCAM Difference Map}
Compute per-class GradCAM maps:
\begin{equation}
M_c = \text{GradCAM}(x, c) \quad \forall c \in \{0,1\}
\end{equation}
Difference map highlights spurious regions:
\begin{equation}
M_{\text{diff}} = |M_{\hat{y}} - M_{y}|
\end{equation}
where $\hat{y}$ is predicted class, $y$ is true label. High values indicate regions where predicted class attribution differs from true class.

\textbf{Step 2: Percentile Thresholding}
Binary mask isolates top-$p$ spurious pixels:
\begin{equation}
\mathcal{S} = \{(i,j) : M_{\text{diff}}(i,j) \geq \text{percentile}(M_{\text{diff}}, p)\}
\end{equation}
Default: $p = 75$ (top 25\% most spurious). \textbf{Note:} This choice is not empirically validated --- grid search over $p \in \{50, 75, 90\}$ deferred to full experiment.

\textbf{Step 3: Gradient Variance Penalty}
Measure gradient variance in spurious region:
\begin{equation}
L_{\text{spatial}} = \text{Var}\left(\nabla_x f_\theta(x) \odot \mathbb{1}_\mathcal{S}\right)
\end{equation}
where $\odot$ is element-wise product, $\mathbb{1}_\mathcal{S}$ is the binary mask. Higher variance in spurious regions indicates scattered gradients.

\textbf{Step 4: Adaptive Lambda Scaling}
Adjust penalty strength based on worst-group accuracy on validation set:
\begin{equation}
\lambda_{t+1} = \begin{cases}
\lambda_t \cdot 1.2 & \text{if } \text{WGA}_{\text{val}} < \text{target} \\
\lambda_t \cdot 0.8 & \text{if } \text{WGA}_{\text{val}} > \text{target} + 5\% \\
\lambda_t & \text{otherwise}
\end{cases}
\end{equation}
Initial: $\lambda_0 = 0.01$. Prevents under/over-regularization.

\subsubsection{Theoretical Justification}

\textbf{Mechanism:} Spurious gradients are dense (high variance) due to conflict between core and spurious features. Penalizing variance in spurious regions forces the model to rely on core features, reducing spurious shortcut.

\textbf{Why percentile normalization?} Absolute thresholding biases toward majority samples (majority gradients dominate distribution). Percentile thresholding selects top-$p$ spurious pixels \emph{per sample}, avoiding majority bias.

\textbf{Why adaptive lambda?} Fixed penalty may be too weak (no WGA improvement) or too strong (catastrophic accuracy drop). Adaptive scaling maintains balance.

\subsubsection{Hyperparameters}

\begin{table}[h]
\centering
\caption{Spatial regularization hyperparameters}
\label{tab:hyperparams}
\begin{tabular}{lll}
\toprule
Parameter & Value & Rationale \\
\midrule
$\lambda_0$ & 0.01 & Conservative initial penalty \\
Percentile $p$ & 75 & Top quartile (MNIST PoC) \\
Scaling factor & 1.2 / 0.8 & Gradual adjustment (20\% steps) \\
WGA target & 75\% & Dataset-dependent \\
\bottomrule
\end{tabular}
\end{table}

\subsection{Implementation Details}

\textbf{Framework:} PyTorch 2.1+ (required for H100 sm\_90 support)
\textbf{GradCAM library:} pytorch-grad-cam (pip install grad-cam)
\textbf{SegFormer:} HuggingFace transformers (semantic segmentation for background swap)
\textbf{Statistical tests:} scipy.stats (Welch's t-test, Pearson correlation)

\textbf{Reproducibility:} All experiments use seed=42. Code released upon publication.

\subsection{Assumptions}

\textbf{A1 (GradCAM Validity):} Minority samples correctly classified $\geq$60\% (else GradCAM may not localize to bird region).
\textbf{A2 (Percentile Normalization):} 75th percentile avoids majority bias (lower percentiles may include non-spurious regions).
\textbf{A3 (Causal Attribution):} Gradient scattering \emph{caused} by spurious conflict, not sample complexity. Validated via background swap test.
\textbf{A4 (Regularization Mechanism):} Spatial penalty reduces spurious reliance, not just smooths gradients. Validated via WGA improvement in PoC.
